\chapter{The Labeled Transition System and Context-Decorated HML}
\label{sec:lts-hml}
% ===================================================================

\section{Minimal Contexts and the LTS}

To obtain a bisimulation that is a \emph{congruence} for a GSLT, we follow
the approach of Milner, Sewell, and Leifer~\cite{leifer2000deriving}.
The key idea is to label transitions not merely by actions but by
\emph{minimal reactive contexts} --- the smallest environments needed to
trigger a given reduction.

\begin{definition}[Context]
A \emph{context} $K[-]$ for a GSLT $\GSLT$ is a term of $\terms(\GSLT)$
with a distinguished \emph{hole} $[-]$.  The application $K[P]$ substitutes
$P$ for the hole.  A context is \emph{minimal} for a term $P$ and rule $r$ if:
\begin{enumerate}[label=(\roman*)]
  \item $K[P] \rewrite_r Q$ for some $Q$;
  \item there is no proper sub-context $K'$ of $K$ satisfying (i).
\end{enumerate}
\end{definition}

The Milner--Sewell--Leifer construction produces, for any GSLT $\GSLT$, a
labeled transition system $\lts(\GSLT)$ in which transitions are labeled by
minimal contexts:
\[
  P \xrightarrow{K} Q \quad\text{iff}\quad K[P] \rewrite Q.
\]
The bisimulation induced by $\lts(\GSLT)$ is a congruence.

\section{Context-Decorated Hennessy--Milner Logic}

\begin{definition}[Context-Decorated HML]
\label{def:cdHML}
The formulae of \emph{context-decorated Hennessy--Milner logic}
$\HML(\ctx)$ are generated by:
\[
  \phi, \psi \;::=\; \top \;\mid\; \bot \;\mid\; \phi \wedge \psi \;\mid\;
  \neg\phi \;\mid\; \langle K \rangle \phi
\]
where $K$ ranges over contexts of $\GSLT$.
The satisfaction relation is:
\[
  u \sat \langle K \rangle \phi \quad\text{iff}\quad
  K[u] \rewrite u' \;\text{ in one step, and }\; u' \sat \phi.
\]
\end{definition}

\begin{theorem}[Adequacy]
For any GSLT $\GSLT$ equipped with the Milner--Sewell--Leifer LTS:
\[
  u \bisim v \quad\Longleftrightarrow\quad
  \forall \phi \in \HML(\ctx),\; u \sat \phi \Leftrightarrow v \sat \phi.
\]
\end{theorem}

\section{The Logical Metric}

Fix an enumeration $(\phi_n)_{n \in \NN}$ of the formulae of $\HML(\ctx)$,
ordered compatibly with the partial order on contexts (smaller contexts
appearing earlier).

\begin{definition}[Logical Metric]
\[
  d_{\HML}(u, v) \;=\; 2^{-n}, \quad\text{where}\quad
  n = \min\bigl\{k \;\mid\; u \sat \phi_k \not\Leftrightarrow v \sat \phi_k\bigr\}.
\]
If $u \bisim v$, set $d_{\HML}(u, v) = 0$.
\end{definition}

\begin{proposition}
$d_{\HML}$ is an \emph{ultrametric} on $\terms(\GSLT)/{\bisim}$.
Terms differing only in behaviors triggered by large (expensive) contexts
are closer under $d_{\HML}$ than terms differing in behaviors triggered
by small (cheap) contexts.  This is the \emph{locality principle}: local
distinguishability weighs more than global.
\end{proposition}

% ===================================================================
